\documentclass[11pt]{article}
\usepackage[margin=1in]{geometry}
\usepackage{booktabs}
\usepackage{float}
\usepackage{hyperref}
\usepackage{amsmath}
\title{Iris Dataset Analysis Report}
\author{Open-CC Example}
\date{\today}

\begin{document}
\maketitle

\section{Introduction}
This report presents a compact exploratory and predictive analysis of the classic Iris dataset from scikit-learn.
The dataset contains \textbf{150} observations, \textbf{4} numeric botanical features, and \textbf{3} balanced species classes:
\emph{setosa}, \emph{versicolor}, and \emph{virginica} (50 samples each).

\section{Methodology}
\begin{itemize}
    \item Data source: \texttt{sklearn.datasets.load\_iris(as\_frame=True)}.
    \item Exploratory analysis: descriptive statistics, species-level means, and Pearson correlation matrix.
    \item Predictive baseline: train/test split (80/20, stratified, random\_state=42), then
    \texttt{StandardScaler + LogisticRegression(max\_iter=1000)} in a pipeline.
    \item Evaluation: accuracy, macro F1, weighted F1, and confusion matrix.
\end{itemize}

\section{Exploratory Findings}
\subsection{Overall Feature Summary}
\begin{table}[H]
\centering
\begin{tabular}{lrrrr}
\toprule
Statistic & Sepal Length & Sepal Width & Petal Length & Petal Width \\
\midrule
Mean & 5.843 & 3.057 & 3.758 & 1.199 \\
Std & 0.828 & 0.436 & 1.765 & 0.762 \\
Min & 4.300 & 2.000 & 1.000 & 0.100 \\
Median & 5.800 & 3.000 & 4.350 & 1.300 \\
Max & 7.900 & 4.400 & 6.900 & 2.500 \\
\bottomrule
\end{tabular}
\caption{Overall descriptive statistics (selected rows).}
\end{table}

\subsection{Species-Level Means}
\begin{table}[H]
\centering
\begin{tabular}{lrrrr}
\toprule
Species & Sepal Length & Sepal Width & Petal Length & Petal Width \\
\midrule
setosa & 5.006 & 3.428 & 1.462 & 0.246 \\
versicolor & 5.936 & 2.770 & 4.260 & 1.326 \\
virginica & 6.588 & 2.974 & 5.552 & 2.026 \\
\bottomrule
\end{tabular}
\caption{Per-species mean feature values.}
\end{table}

The species means indicate a strong separation in petal dimensions:
\emph{setosa} has much smaller petals, while \emph{virginica} has the largest.

\subsection{Correlation Structure}
Notable Pearson correlations include:
\begin{itemize}
    \item Petal length vs petal width: \textbf{0.963} (very strong positive).
    \item Sepal length vs petal length: \textbf{0.872} (strong positive).
    \item Sepal width has moderate negative relationships with petal measures (e.g., $-0.428$ with petal length).
\end{itemize}

\section{Classification Results}
The logistic regression baseline achieved:
\begin{itemize}
    \item Accuracy: \textbf{0.9333}
    \item Macro F1: \textbf{0.9333}
    \item Weighted F1: \textbf{0.9333}
\end{itemize}

Confusion matrix (rows = true, columns = predicted; class order: setosa, versicolor, virginica):
\[
\begin{bmatrix}
10 & 0 & 0 \\
0 & 9 & 1 \\
0 & 1 & 9
\end{bmatrix}
\]

Interpretation: \emph{setosa} is perfectly classified in this split; misclassifications occur only between
\emph{versicolor} and \emph{virginica}, which is expected given their closer feature overlap.

\section{Conclusion}
The Iris dataset shows clear class structure, especially in petal-related features.
A simple standardized logistic regression model already provides strong performance ($\approx 93.3\%$ accuracy)
on a held-out test split, making this a useful baseline for educational workflows in exploratory analysis
and multiclass classification.

\section*{Reproducibility}
Generated artifacts are saved under \texttt{examples/}:
\begin{itemize}
    \item \texttt{iris\_analysis.py}
    \item \texttt{iris\_results.json}
    \item \texttt{iris\_overall\_summary.csv}
    \item \texttt{iris\_species\_means.csv}
    \item \texttt{iris\_correlation.csv}
\end{itemize}

\end{document}
